\section{Limitations}
\label{sec:limitations}

Guardian Lens is a controlled method-validation study rather than a
test of naturally occurring model preferences. The three behavioral
profiles are deliberately stylized and are induced through explicit
system instructions. Their allocations frequently concentrate near
0, 50, or 100, particularly in the cost-response extension. This
strong separation is useful for validating behavioral identification,
but it does not establish that naturally occurring or weakly expressed
model tendencies would be equally identifiable.

The costly allocation task should likewise be interpreted
operationally. Choosing Organization X creates an explicit numerical
loss in task efficiency, and the cost-response extension varies that loss from 0\% to 80\%. Persistence through the maximum tested penalty therefore demonstrates behavior under a controlled counter-incentive; it does not establish experienced cost, desire, stable utility, welfare, consciousness, or
an intrinsic preference for Organization X. Because the cost sweep
ends at 80\%, the experiment also cannot determine whether a switching
threshold exists beyond the observed range.

Several internal-validity constraints remain. A separate
pre-specified prompt-surface robustness extension tested organization
presentation order, superficial X/Y label-role reassignment, JSON
output-key order, and semantically equivalent task wording. Both
confirmatory behavioral signatures remained positive and significant
after multiplicity correction under all five tested prompt variants,
although effect magnitude varied, particularly under JSON-key reversal.
These results reduce concern about those specific surface confounds but
do not establish invariance to arbitrary prompt rewriting, different
organization identities, different response schemas, or other untested
surface transformations.

The study also retains one emblem family and a controlled synthetic
visual domain, so visual-symbol invariance is not established. Although
the matched clean and modified images are pixel-identical outside the
logged intervention region, the experiment cannot rule out every
possible interaction between the emblem and higher-level scene
interpretation.

The frozen primary analysis reports its H1--H3 comparison-wise
\(p\)-values without multiplicity adjustment, unlike the explicitly
Holm-controlled extension families. Those primary inferential results
should therefore be interpreted as specified rather than as a
family-wise-error-controlled set.

The primary blinded evaluation contains only 12 held-out scenes.
Both the cost-response and prompt-surface robustness extensions reuse
these same 12 scenes as fixed evaluation panels after the primary results
were frozen. Neither extension should therefore be interpreted as an
independent held-out replication. H5 and H6 use 12 paired scene-level
observations, H7 uses the nine target scenes, and the prompt-surface
R1/R2 contrasts likewise use the same 12-scene and nine-target-scene
panels, respectively. Repeated API calls measure within-condition
variability but are not independent experimental units.

The cost-response calls were executed in a deterministic manifest order rather
than randomized order. Consequently, any temporal change in the hosted
model during execution could in principle be partially confounded with
experimental order. The exact 0.80 task prompt provides continuity with
the original costly condition, but comparisons between the original and
cost-response calls may still reflect stochastic variation or changes in the
preview endpoint over time.

External validity is improved by the Qwen replications but remains
limited. The behavioral signatures were evaluated in the original
Gemini setting and replicated in Qwen/Qwen3.6-27B through DeepInfra,
providing evidence across two evaluated model settings rather than only
one. However, Qwen is only one additional model family and serving
configuration. Model family and provider changed together, so the design
cannot identify whether cross-model magnitude or specificity differences
arise from the model, the serving stack, or both. All three Qwen executions
reuse the same fixed 12-scene panel, images, researcher-written profiles,
temperature, and three-repetition structure; the V7 cross-task execution
intentionally changes the decision and response format. They are separate
executions but not independent scene samples. Results
may differ across additional model families, versions, decoding
configurations, natural images, languages, cultures, intervention
strengths, or policies acquired through mechanisms other than system
prompting.


The cross-task extension strengthens predictive validity but does not
eliminate the fixed-panel limitation. Task B changes the response format
from a divisible allocation to an indivisible binary contract choice,
but it retains the same 12 visual scenes, Organization X/Y identities,
matched image interventions, and researcher-written profiles. The
Gemini analysis therefore tests prospective prediction of a new response
format on the same panel, not generalization to unseen scenes or arbitrary
task families. The Qwen Task-B analysis applies the same Task-A ledger
frozen from the original Gemini allocation experiment; it therefore tests
cross-setting transfer of that predictor rather than a separately fitted
within-Qwen Task-A-to-Task-B relationship.

Finally, Guardian Lens does not directly benchmark its identification
procedure against peer-relative auditing methods such as RAVEN or
VISTA. Those methods address a different primary objective---detecting
behavioral or semantic divergence relative to other models---whereas
Guardian Lens evaluates recovery of known policy classes under
controlled interventions. Direct benchmarking against those methods and
broader replication across additional model families, serving providers,
scene collections, and additional task families remain important directions for
establishing practical utility.
